%% toc: P2 21 21.1,21.8
\section{Validity, finite satisfiability and semi-conservative reductions}
\label{sec:p02-w02b-p02-031-validity-finsat}

\emph{Purpose.}  \Cref{sec:p02-w02-p02-036-finite-validity} sorts the
first-order formulas into the finitely satisfiable ones, the infinity axioms
and the unsatisfiable ones, and proves that the three classes are pairwise
recursively inseparable.  This section adds further results from the same
development of Börger, Grädel and Gurevich
\cite[Sect.~2.1.2, pp.~36--38]{BorgerGradelGurevich1997}.  The first group
proves the same inseparability with formulas that describe two-register
machines, and uses those formulas for one reduction that preserves
satisfiability and finite satisfiability together, with values in two prefix
classes of Krom and Horn sentences.  The next result shows
that unrestricted validity, although it is a different question from finite
satisfiability, can be decided by anyone who can decide finite
satisfiability: there is a computable translation that sends the valid
formulas, and only those, to finitely satisfiable sentences.  The last group
introduces the semi-conservative reductions, which are required to respect
only the unsatisfiable and the finitely satisfiable formulas, and shows what
such a reduction does and does not carry over.

\emph{Inputs.}  The fixed effective coding of first-order formulas of
\cref{sec:decision-hierarchy-01}; the countermodel classes of
Definition~\ref{def:p02-w02-p02-036-classes}, recursive inseparability
(Definition~\ref{def:p02-w02-p02-036-inseparable}) and conservative
reductions (Definition~\ref{def:p02-w02-p02-036-conservative}); the run
sentence \(\psi_{\mathcal T,w}\) of
Definition~\ref{def:p02-w02-p02-036-run-sentence}, the path translation
\(\varphi\mapsto\varphi^{\sharp}\) of
Definition~\ref{def:p02-w02-p02-036-gadget}, and the renaming
\(\chi\mapsto\chi[R/E]\) of Lemma~\ref{thm:p02-w02-p02-036-transfer}.  A formula
is \emph{valid} when it is satisfied in every structure for its signature
under every assignment.  As in the source, ``recursive'' means
``decidable'' for sets and ``total computable'' for functions.

\emph{Conclusions.}  Two-register machines simulate Turing machines with the
halting state recorded (\cref{thm:p02-w02b-p02-031-two-register-simulation});
the register-machine formulas reduce the halting problems to unsatisfiability,
finite satisfiability and infinity axioms
(\cref{thm:p02-w02b-p02-031-register-reduction}); this gives Trakhtenbrot's
theorem again (\cref{thm:p02-w02b-p02-031-ex-2-1-32}) and a conservative
reduction of all formulas into each of the classes
\([\exists\forall\exists\forall]\cap\mathrm{KROM}\cap\mathrm{HORN}\) and
\([\forall\exists\exists\forall]\cap\mathrm{KROM}\cap\mathrm{HORN}\) with only
binary predicate symbols
(\cref{thm:p02-w02b-p02-031-ex-2-1-34,thm:p02-w02b-p02-031-ex-2-1-37}).
For every finite relational signature \(\Lang\) with a
binary relation symbol there is a computable map from first-order formulas
to \(\Lang\)-sentences that sends exactly the valid formulas into
\(\mathrm{Fin\text{-}sat}(\operatorname{Sent}(\Lang))\)
(\cref{thm:p02-w02b-p02-031-val-to-finsat}).  A semi-conservative reduction
carries the recursive inseparability of the unsatisfiable and the finitely
satisfiable formulas from its domain to its target
(\cref{thm:p02-w02b-p02-031-semi-inseparable}).  A semi-conservative
reduction between arbitrary recursive classes need not be accompanied by a
conservative one (\cref{thm:p02-w02b-p02-031-semi-not-conservative}), but
one out of the class \(\mathrm{FO}\) of all formulas always is
(\cref{thm:p02-w02b-p02-031-gurevich}); the proof runs through effective
inseparability and an embedding theorem of recursion theory
(\cref{thm:p02-w02b-p02-031-eff-classes,thm:p02-w02b-p02-031-embedding}).
All proofs are in the annex.

\subsection{Register-machine formulas and Krom--Horn conservative reduction classes}
\label{sec:p02-w02b-p02-031-register}

The source gives a second proof of Trakhtenbrot's theorem, in which the
formulas describe two-register machines instead of Turing machines
(its Exercise~2.1.32).  Applied to one fixed machine that searches at the same
time for a refutation and for a finite model of a given formula, the same
formulas give a single reduction that respects satisfiability and finite
satisfiability together (Exercise~2.1.34), and its values lie in two prefix
classes of Krom and Horn sentences (Exercise~2.1.37)
\cite[Sect.~2.1.2, pp.~36--37]{BorgerGradelGurevich1997}.  The machines and the
shape of the formulas are those of the source's Theorem~2.1.15 and
Corollary~2.1.16 \cite[Sect.~2.1.1, pp.~27--29]{BorgerGradelGurevich1997}.

\emph{Purpose and inputs.}  The subsection needs the Turing machines of
\cref{sec:turing-machines-01}, the machine codes \(V\), \(M_e\) and the halting
problems \(H_1\), \(H_2\), \(H_\infty\) of
Definition~\ref{def:p02-w02-p02-036-halting}, the classes of
Definition~\ref{def:p02-w02-p02-036-classes}, and the index constructions of
\cref{thm:p02-w02b-p02-031-code-constructions}.  Its conclusions are a
reduction property for the register-machine formulas
(\cref{thm:p02-w02b-p02-031-register-reduction}), and from it
Trakhtenbrot's theorem in the source's register-machine form
(\cref{thm:p02-w02b-p02-031-ex-2-1-32}), the simultaneous reduction
(\cref{thm:p02-w02b-p02-031-ex-2-1-34}) and the two conservative reduction
classes (\cref{thm:p02-w02b-p02-031-ex-2-1-37}).

\begin{hbdefinition}{def:p02-w02b-p02-031-two-register}{Two-register machines}
A \emph{two-register machine program} \(M\) has a finite set
\(Q_M\subseteq\mathbb N\) of states with \(0,1,2\in Q_M\).  The states \(1\)
and \(2\) are the \emph{stop states} and carry no instruction; every other
state \(i\in Q_M\) carries exactly one instruction \(I_i\), of one of two
kinds \cite[Sect.~2.1.1, pp.~27--29]{BorgerGradelGurevich1997}:
\begin{itemize}
\item an \emph{addition instruction} \(I_i=(i,r,j)\) with \(r\in\{1,2\}\) and
      \(j\in Q_M\): at state \(i\), add \(1\) to register \(r\) and go to
      state \(j\);
\item a \emph{subtraction instruction} \(I_i=(i,r,j,k)\) with \(r\in\{1,2\}\)
      and \(j,k\in Q_M\): at state \(i\), if register \(r\) contains \(0\), go
      to state \(j\); otherwise subtract \(1\) from register \(r\) and go to
      state \(k\).
\end{itemize}
A \emph{configuration} is a triple \((i,m,n)\in Q_M\times\mathbb N\times
\mathbb N\): state \(i\), content \(m\) of register \(1\) and content \(n\) of
register \(2\).  The instruction of a state other than \(1\) and \(2\)
determines a unique next configuration; a configuration in a stop state has
none.  The \emph{computation} of \(M\) from a configuration \(C\) is the unique
maximal finite or infinite sequence of configurations that begins with \(C\)
and proceeds by next configurations.  It \emph{halts} if it is finite, that
is, if it reaches a stop state.  Write \(C\Rightarrow_M D\) if \(D\) occurs in
the computation from \(C\).  The \emph{register contents} of a computation are
the numbers \(m\) and \(n\) of its configurations \((i,m,n)\).
\index{register machine!two-register}\index{two-register machine}
\end{hbdefinition}

The source allows halting states in general and restricts attention to
programs with the two stop states \(1\) and \(2\) for this application
\cite[Sect.~2.1.2, p.~34]{BorgerGradelGurevich1997}; the definition builds
that restriction in.

\begin{hbdefinition}{def:p02-w02b-p02-031-krom-horn}{Krom and Horn formulas; prefix classes}
A \emph{clause} is a disjunction of atomic and negated atomic formulas.  A
\emph{Horn clause} is a clause with at most one non-negated atom; it can be
written as an implication \(\alpha\to\beta\) in which \(\alpha\) is a
conjunction of atoms and \(\beta\) is an atom or the constant
\(\mathit{false}\).  A \emph{Krom clause} is a clause with at most two
constituents.  \(\mathrm{HORN}\) (respectively \(\mathrm{KROM}\)) is the
class of first-order formulas in prenex normal form whose quantifier-free part
is a conjunction of Horn (respectively Krom) clauses
\cite[Sect.~2.1.1, Def.~2.1.14, p.~27]{BorgerGradelGurevich1997}.  For a
string \(w\) of quantifier symbols, \([w]\) is the class of prenex sentences of
first-order logic without equality and without function symbols whose
quantifier prefix is \(w\), as in the prefix-vocabulary classes of
\cite[Sect.~1.3, Def.~1.3.1, p.~8]{BorgerGradelGurevich1997}.  The phrase
``with only binary predicate symbols'' restricts a class to sentences all of
whose predicate symbols are binary.
\index{Krom formula}\index{Horn formula}\index{prefix class}
\end{hbdefinition}

\begin{hbtheorem}{thm:p02-w02b-p02-031-two-register-simulation}{proved=annex:p02-w02b-p02-031-two-register-simulation}{Two-register machines simulate Turing machines}
For every Turing machine \(\mathcal T\) there are a two-register machine
program \(M_{\mathcal T}\) and a map \(w\mapsto m_{\mathcal T}(w)\) on input
words, both computable from the transition table of \(\mathcal T\), such that
for every input word \(w\) the computation of \(M_{\mathcal T}\) from
\((0,m_{\mathcal T}(w),0)\) has the following properties.
\begin{enumerate}
\item It reaches \((1,0,0)\) if and only if the computation of \(\mathcal T\)
      on \(w\) reaches \(q_{\mathrm{rej}}\).
\item It reaches \((2,0,0)\) if and only if the computation of \(\mathcal T\)
      on \(w\) reaches \(q_{\mathrm{acc}}\).
\item Every configuration of it in a stop state is \((1,0,0)\) or \((2,0,0)\).
\item If the computation of \(\mathcal T\) on \(w\) never halts, then it never
      halts and its register contents are unbounded.
\end{enumerate}
Moreover \(m_{\mathcal T}(\varepsilon)=1\) for the empty word
\(\varepsilon\), and there is a program \(M^{0}_{\mathcal T}\), computable from
the table of \(\mathcal T\), whose computation from \((0,0,0)\) has
properties 1--4 for \(w=\varepsilon\).
\end{hbtheorem}

The source takes the universality of two-register machines from the
literature \cite[Sect.~2.1.1, p.~28]{BorgerGradelGurevich1997}; the annex proves
the form needed here, in which the stop state records the halting state of
\(\mathcal T\), the registers are emptied before stopping, and a
non-halting computation drives the registers to infinity.  The last property
is the register-machine counterpart of the source's sharpened stop condition,
by which a computation that does not halt never becomes cyclic
\cite[Sect.~2.1.2, p.~35]{BorgerGradelGurevich1997}.

\begin{hbdefinition}{def:p02-w02b-p02-031-register-formulas}{The register-machine reduction formulas}
Let \(M\) be a two-register machine program and \(m\in\mathbb N\).  The
signature consists of binary predicate symbols \(K\), \(K_i\) for \(i\in
Q_M\), and, when \(m\geq1\), \(A_0,\ldots,A_m\).  A configuration
\((i,a,b)\) is described by the atom \(\underline{(i,a,b)}\defeq K_iab\), the
register contents being read as elements
\cite[Sect.~2.1.1, p.~28]{BorgerGradelGurevich1997}.  Let
\(\alpha_{M,m}(u,x,v,y)\) be the conjunction of the following clauses, in
which \(u\) plays the part of the number \(0\) and \(v\) that of the successor
of \(x\).
\begin{enumerate}
\item \(\mathrm{STEP}_M\): for each addition instruction \((i,1,j)\) the clause
      \(K_ixy\to K_jvy\), and for each \((i,2,j)\) the clause
      \(K_iyx\to K_jyv\); for each subtraction instruction \((i,1,j,k)\) the
      clauses \(K_ivy\to K_kxy\) and \(K_iuy\to K_juy\), and for each
      \((i,2,j,k)\) the clauses \(K_iyv\to K_kyx\) and \(K_iyu\to K_jyu\).
\item \(\mathrm{START}_m\): for \(m=0\) the atom \(K_0uu\); for \(m\geq1\) the
      atom \(A_0uu\), the clauses \(A_lxy\to A_{l+1}vy\) for \(l<m\), and the
      clause \(A_mxy\to K_0xy\).
\item \(\mathrm{NOSTOP}\): the clause \(\neg K_1uu\).
\item \(\prec_{K,M}\): the clauses \(\neg Kxx\) and \(Kyx\to Kyv\), and for
      each \(i\in Q_M\) the clauses \(K_ixy\to Kxv\) and \(K_iyx\to Kxv\).
\end{enumerate}
Put
\[
 \beta_{M,\prec,m}\defeq\exists u\,\forall x\,\exists v\,\forall y\;
   \alpha_{M,m},
 \qquad
 \beta^{+}_{M,\prec,m}\defeq\forall x\,\exists u\,\exists v\,\forall y\;
   \alpha_{M,m},
\]
and \(\psi_{M,\prec}\defeq\beta_{M,\prec,0}\).  In the source's notation
\(\beta_{M,\prec,m}\) is \(\beta_{M,\prec,C,D}\) with \(C=(0,m,0)\) and
\(D=(1,0,0)\) \cite[Sect.~2.1.2, p.~36]{BorgerGradelGurevich1997}.
\index{reduction formula!register machine}
\end{hbdefinition}

Eliminating the existential quantifiers from left to right, as in the source's
Skolemization \cite[Sect.~2.1.1, Thm.~2.1.8 and Lemma~2.1.9,
pp.~25--26]{BorgerGradelGurevich1997}, replaces \(u\) by a new constant \(0\)
and \(v\) by a term \(x'\) in \(\beta_{M,\prec,m}\), and \(u\) by a term
\(x^{+}\) and \(v\) by \(x'\) in \(\beta^{+}_{M,\prec,m}\).  So
\(\psi_{M,\prec}\) has the Skolem normal form
\[
 \forall x\,\forall y\,\bigl(\mathrm{STEP}_M\land\underline{(0,0,0)}\land
 \neg\underline{(1,0,0)}\land\prec_{K,M}\bigr)
\]
displayed in the source's Exercise~2.1.32, with
\(\prec_{K,M}=\neg Kxx\land(Kyx\to Kyx')\land\bigwedge_i\bigl((K_ixy\to
Kxx')\land(K_iyx\to Kxx')\bigr)\)
\cite[Sect.~2.1.2, p.~36]{BorgerGradelGurevich1997}.  The source prints
\(u\) in place of \(x'\) in this display; in a Skolem normal form with the
variables \(x,y\) only, the letter can only stand for the successor term, as in
the source's relativized axiom \(\forall x\,\exists u\,\forall y(\ldots)\) on
p.~35, where \(u\) is the witness that follows \(x\).  For \(m\geq1\) the
source writes the start atom \(\underline{(0,m,0)}\) with the term
\(0^{(m)}\) for \(m\); a prenex sentence with prefix \(\exists\forall\exists\forall\)
has no such term, and the clauses with \(A_0,\ldots,A_m\) derive that atom
instead.

Every conjunct of \(\alpha_{M,m}\) is a Krom clause and a Horn clause, all
predicate symbols are binary, and there are no function symbols and no
equality.  So
\(\beta_{M,\prec,m}\in[\exists\forall\exists\forall]\cap\mathrm{KROM}\cap
\mathrm{HORN}\) and \(\beta^{+}_{M,\prec,m}\in[\forall\exists\exists\forall]
\cap\mathrm{KROM}\cap\mathrm{HORN}\), both with only binary predicate
symbols.  Each conjunct mentions either none of \(u\) or none of \(x,v\);
this is what lets one matrix serve both prefixes.

\begin{hbtheorem}{thm:p02-w02b-p02-031-register-reduction}{proved=annex:p02-w02b-p02-031-register-reduction}{Reduction property of the register-machine formulas}
Let \(M\) be a two-register machine program and \(m\in\mathbb N\), and assume
that the computation of \(M\) from \((0,m,0)\) satisfies: every configuration
of it in a stop state is \((1,0,0)\) or \((2,0,0)\); and if it never halts,
its register contents are unbounded.  Let \(\beta\) be
\(\beta_{M,\prec,m}\) or \(\beta^{+}_{M,\prec,m}\).  Then:
\begin{enumerate}
\item \((0,m,0)\Rightarrow_M(1,0,0)\) if and only if \(\beta\) is
      contradictory;
\item \((0,m,0)\Rightarrow_M(2,0,0)\) if and only if \(\beta\) is finitely
      satisfiable; in that case \(\beta\) has a model with \(N+2\) elements,
      where \(N\) is the largest register content of the computation;
\item the computation never halts if and only if \(\beta\) is an infinity
      axiom.
\end{enumerate}
\end{hbtheorem}

For \(m=0\), items 1 and~2 are the source's reduction property for
\(\psi_{M,\prec}\) with the halting problems of register machines in place of
those of Turing machines \cite[Sect.~2.1.2, p.~34]{BorgerGradelGurevich1997}.

\begin{hbcorollary}{thm:p02-w02b-p02-031-ex-2-1-32}{proved=annex:p02-w02b-p02-031-ex-2-1-32}{Trakhtenbrot's theorem with register-machine formulas}
For \(e\in V\) let \(\psi_e\defeq\psi_{M^{0}_{M_e},\prec}\), with
\(M^{0}_{M_e}\) as in \cref{thm:p02-w02b-p02-031-two-register-simulation}.
Then \(e\in H_1\), \(e\in H_2\) and \(e\in H_\infty\) hold according as
\(\psi_e\) is contradictory, finitely satisfiable or an infinity axiom.
Consequently, for every recursive class \(Y\) of formulas that contains every
\(\psi_e\), in particular for \(Y=\mathrm{FO}\) and for
\(Y=[\exists\forall\exists\forall]\cap\mathrm{KROM}\cap\mathrm{HORN}\) with
only binary predicate symbols, the sets \(\mathrm{Fin\text{-}sat}(Y)\),
\(\mathrm{Non\text{-}sat}(Y)\) and \(\mathrm{Inf\text{-}axioms}(Y)\) are
pairwise recursively inseparable
\cite[Sect.~2.1.2, Thm.~2.1.30 and Exercise~2.1.32]{BorgerGradelGurevich1997}.
\end{hbcorollary}

The source's logic machine of Corollary~2.1.16 recognizes the unsatisfiable
formulas.  For the simultaneous reduction it is enriched so that it also
recognizes the finitely satisfiable ones, with the two halting problems
\(H'_{M,i}=\{\psi:C(\psi)\Rightarrow_M(i,0,0)\}\) for \(i=1,2\), where
\(C(\psi)=(0,m(\psi),0)\) \cite[Sect.~2.1.2, p.~36]{BorgerGradelGurevich1997}.

\begin{hbtheorem}{thm:p02-w02b-p02-031-logic-machine}{proved=annex:p02-w02b-p02-031-logic-machine}{A logic machine for unsatisfiability and finite satisfiability}
There are a two-register machine program \(M_L\) and a computable map
\(\psi\mapsto m(\psi)\) on \(\mathrm{FO}\) such that for every first-order
formula \(\psi\) the computation of \(M_L\) from \(C(\psi)\defeq(0,m(\psi),0)\)
reaches \((1,0,0)\) if \(\psi\in\mathrm{Non\text{-}sat}\), reaches \((2,0,0)\)
if \(\psi\in\mathrm{Fin\text{-}sat}\), and never halts, with unbounded register
contents, if \(\psi\in\mathrm{Inf\text{-}axioms}\); every configuration of it
in a stop state is \((1,0,0)\) or \((2,0,0)\).  So
\(H'_{M_L,1}=\mathrm{Non\text{-}sat}\) and \(H'_{M_L,2}=\mathrm{Fin\text{-}sat}\).
\end{hbtheorem}

\begin{hbcorollary}{thm:p02-w02b-p02-031-ex-2-1-34}{proved=annex:p02-w02b-p02-031-ex-2-1-34}{The refined logic-machine formulas reduce both problems at once}
With \(M_L\) and \(m\) as in \cref{thm:p02-w02b-p02-031-logic-machine}, put
\(g(\psi)\defeq\beta_{M_L,\prec,m(\psi)}\), whose Skolem normal form is
\(\forall x\,\forall y\,(\mathrm{STEP}_{M_L}\land\mathrm{START}_{m(\psi)}\land
\neg\underline{(1,0,0)}\land\prec_{K,M_L})\).  For every first-order formula
\(\psi\):
\begin{enumerate}
\item \(\psi\) is satisfiable if and only if \(g(\psi)\) is satisfiable;
\item \(\psi\) is finitely satisfiable if and only if \(g(\psi)\) is finitely
      satisfiable;
\item more precisely, \(g\) maps \(\mathrm{Non\text{-}sat}\),
      \(\mathrm{Fin\text{-}sat}\) and \(\mathrm{Inf\text{-}axioms}\) into the
      classes of the same names.
\end{enumerate}
The map \(g\) is computable \cite[Sect.~2.1.2, Exercise~2.1.34]{BorgerGradelGurevich1997}.
\end{hbcorollary}

\begin{hbcorollary}{thm:p02-w02b-p02-031-ex-2-1-37}{proved=annex:p02-w02b-p02-031-ex-2-1-37}{Two conservative reduction classes of Krom and Horn sentences}
The classes \([\exists\forall\exists\forall]\cap\mathrm{KROM}\cap
\mathrm{HORN}\) and \([\forall\exists\exists\forall]\cap\mathrm{KROM}\cap
\mathrm{HORN}\) with only binary predicate symbols are conservative reduction
classes in the sense of Definition~\ref{def:p02-w02-p02-036-conservative}:
\(g\) of \cref{thm:p02-w02b-p02-031-ex-2-1-34} is a conservative reduction from
\(\mathrm{FO}\) to the first, and
\(g^{+}(\psi)\defeq\beta^{+}_{M_L,\prec,m(\psi)}\) is one to the second
\cite[Sect.~2.1.2, Exercise~2.1.37]{BorgerGradelGurevich1997}.
\end{hbcorollary}

By \cref{thm:p02-w02-p02-036-conservative}, the sets
\(\mathrm{Non\text{-}sat}(X)\), \(\mathrm{Fin\text{-}sat}(X)\) and
\(\mathrm{Inf\text{-}axioms}(X)\) are pairwise recursively inseparable for
each of these two classes \(X\).  For the second class, the source's
Exercise~2.1.17 suggests the same device in the plain reduction-class setting:
the term \(x^{+}\) serves as the number \(0\)
\cite[Sect.~2.1.1, p.~29]{BorgerGradelGurevich1997}.

\begin{hbexample}{def:p02-w02b-p02-031-register-example}{Three small register machines}
Let \(M\) have the states \(0,1,2,3\) with \(I_0=(0,1,3)\) and
\(I_3=(3,1,2,3)\).  From \((0,0,0)\) it computes \((0,0,0)\), \((3,1,0)\),
\((3,0,0)\), \((2,0,0)\), and halts in the stop state \(2\) with empty
registers; the largest register content is \(N=1\).  The matrix of
\(\psi_{M,\prec}\) consists of \(K_0xy\to K_3vy\), \(K_3vy\to K_3xy\),
\(K_3uy\to K_2uy\), \(K_0uu\), \(\neg K_1uu\) and the clauses of
\(\prec_{K,M}\).  The model built in the annex has \(N+2=3\) elements
\(0,1,2\), with \(u\) interpreted by \(0\), the witness for \(v\) at \(x\)
equal to \(x+1\) for \(x\leq1\) and to \(2\) for \(x=2\),
\(K=\{(0,1),(0,2),(1,2)\}\), \(K_0=\{(0,0)\}\), \(K_3=\{(1,0),(0,0)\}\),
\(K_2=\{(0,0)\}\) and \(K_1=\varnothing\).  For instance, the clause
\(K_3vy\to K_3xy\) has a true premise only at \(x=0\), \(y=0\), where the
witness for \(v\) is \(1\) and \(K_3(0,0)\) holds.

If instead \(I_3=(3,1,1,3)\), the computation ends in \((1,0,0)\), and
\(\psi_{M,\prec}\) is contradictory: in any model, \(K_0uu\) and the step
clauses give \(K_3\) at the pair (witness of \(v\) at \(u\), \(u\)), then
\(K_3uu\), then \(K_1uu\), against \(\neg K_1uu\).  If instead
\(I_3=(3,1,3)\), register \(1\) grows forever, and \(\psi_{M,\prec}\) is an
infinity axiom by \cref{thm:p02-w02b-p02-031-register-reduction}(3).
\end{hbexample}

\begin{hbwarning}{def:p02-w02b-p02-031-register-limits}{What the hypotheses and the formulas do not give}
\begin{enumerate}
\item The hypothesis of \cref{thm:p02-w02b-p02-031-register-reduction} on
  unbounded registers cannot be dropped.  Let \(M\) have
  \(I_0=(0,1,3,3)\) and \(I_3=(3,1,3,3)\).  From \((0,0,0)\) it moves to
  \((3,0,0)\) and stays there for ever, so it never halts.  Yet
  \(\psi_{M,\prec}\) has a two-element model: domain \(\{0,1\}\), \(u\)
  interpreted by \(0\), the witness for \(v\) equal to \(1\) at both
  elements, \(K=\{(0,1)\}\), \(K_0=K_3=\{(0,0)\}\) and \(K_1=K_2=\varnothing\).
  This is why the source sharpens the stop condition before relativizing the
  infinity axiom \cite[Sect.~2.1.2, p.~35]{BorgerGradelGurevich1997}, and why
  \cref{thm:p02-w02b-p02-031-two-register-simulation}(4) is part of the
  simulation.
\item The formulas are frugal: a model may make \(K_i\) true at pairs that the
  computation never visits.  The proofs only use that every visited
  configuration is true in every model.
\item The reduction \(g\) is computable, but \(m(\psi)\) is a power of \(3\)
  whose exponent is a numerical code of the binary numeral of the code of
  \(\psi\), and \(g(\psi)\) has \(m(\psi)+1\) predicate symbols \(A_l\).  The
  construction gives no useful bound on the size of \(g(\psi)\) or of its
  finite models.
\end{enumerate}
\end{hbwarning}

\subsection{An effective reduction of validity to finite satisfiability}
\label{sec:p02-w02b-p02-031-val-to-finsat}

The source states the result as its Exercise~2.1.33
\cite[Sect.~2.1.2, p.~36]{BorgerGradelGurevich1997}, with the hint to use
Herbrand's theorem.  The proof given here takes a shorter route through the
run sentence: the valid formulas form a computably enumerable set, a Turing
machine recognizes that set, and acceptance by a fixed machine is exactly what
the run sentence turns into finite satisfiability.

\begin{hbtheorem}{thm:p02-w02b-p02-031-val-to-finsat}{proved=annex:p02-w02b-p02-031-val-to-finsat}{Validity reduces to finite satisfiability}
Let \(\Lang\) be a finite relational signature containing a binary relation
symbol.  There is a computable map \(F\) from the set \(\mathrm{FO}\) of codes
of first-order formulas to \(\operatorname{Sent}(\Lang)\) such that, for
every first-order formula \(\varphi\),
\[
 \varphi\in\mathrm{Val}
 \quad\Longleftrightarrow\quad
 F(\varphi)\in\mathrm{Fin\text{-}sat}(\operatorname{Sent}(\Lang)).
\]
Consequently a decision procedure for finite satisfiability of
\(\Lang\)-sentences would yield one for validity of first-order formulas
\cite[Sect.~2.1.2, Exercise~2.1.33]{BorgerGradelGurevich1997}.
\end{hbtheorem}

The map \(F\) is a many-one reduction of \(\mathrm{Val}\) to
\(\mathrm{Fin\text{-}sat}(\operatorname{Sent}(\Lang))\) in the sense of
\cref{sec:universal-reductions-02}, once the codes outside \(\mathrm{FO}\) are
sent to a fixed unsatisfiable sentence.  It controls only the finitely
satisfiable side.  A formula that is not valid is sent to a sentence without
a finite model, and the reduction does not say whether that sentence is
unsatisfiable or an infinity axiom.

\begin{hbexample}{def:p02-w02b-p02-031-val-example}{Reading the reduction on two formulas}
Fix \(\Lang=\{R\}\) with \(R\) binary, and let \(\mathcal R\) be the
recognizer of valid sentences used in the proof.  For the valid formula
\(\varphi_1\defeq\forall x\,(Rxx\to Rxx)\), the universal closure is
\(\varphi_1\) itself; \(\mathcal R\) accepts its code after finitely many
steps, say \(s\), so by Lemma~\ref{thm:p02-w02-p02-036-run-sentence} the run
sentence \(\psi_{\mathcal R,\operatorname{bin}(\ulcorner\varphi_1\urcorner)}\)
has a model with \(2s+\lvert\operatorname{bin}(\ulcorner\varphi_1\urcorner)
\rvert+1\) elements, and \(F(\varphi_1)\) has a finite model.  For the formula
\(\varphi_2\defeq Rxy\), which is not valid, the universal closure
\(\forall x\,\forall y\,Rxy\) is false in the two-element structure with
empty \(R\); so \(\mathcal R\) never accepts, the run sentence has no finite
model, and \(F(\varphi_2)\notin\mathrm{Fin\text{-}sat}\).  Nothing in the
construction decides, for \(\varphi_2\), whether \(F(\varphi_2)\) has an
infinite model: that depends on whether the computation of \(\mathcal R\) on
the code rejects or runs forever.
\end{hbexample}

\begin{hbwarning}{def:p02-w02b-p02-031-val-limits}{The reduction runs in one direction only}
\begin{enumerate}
\item No computable \(F'\) satisfies
  \(\varphi\in\mathrm{Val}\iff F'(\varphi)\) is finitely valid.  The finitely
  valid sentences of \(\Lang\) form a co-computably enumerable set
  (Lemma~\ref{thm:p02-w02-p02-036-finsat-ce}(4)); the preimage of a
  co-computably enumerable set under a computable map is co-computably
  enumerable, so \(\mathrm{Val}\) would be co-computably enumerable as well as
  computably enumerable, hence decidable, against
  \cref{thm:fol-validity-undecidable}.
\item \Cref{thm:p02-w02b-p02-031-val-to-finsat} does not give a reduction of
  finite satisfiability to validity, and it is not a conservative reduction:
  it does not preserve satisfiability.
\item The theorem gives no bound on the size of the finite model of
  \(F(\varphi)\) for valid \(\varphi\); that size grows with the number of
  steps the recognizer needs to accept.
\end{enumerate}
\end{hbwarning}

\subsection{Semi-conservative reductions}
\label{sec:p02-w02b-p02-031-semi}

A conservative reduction must respect satisfiability and finite
satisfiability in both directions.  The source observes that, for the
purposes of reduction classes, it can be enough to map the unsatisfiable
formulas to unsatisfiable ones and the finitely satisfiable formulas to
finitely satisfiable ones; nothing is required of the infinity axioms
\cite[Sect.~2.1.2, p.~37]{BorgerGradelGurevich1997}.

\begin{hbdefinition}{def:p02-w02b-p02-031-semi-conservative}{Semi-conservative reduction}
Let \(X\) and \(Y\) be recursive classes of formulas.  A
\emph{semi-conservative reduction} from \(X\) to \(Y\) is a recursive function
\(f:X\to Y\) such that
\begin{enumerate}
\item \(f(\mathrm{Non\text{-}sat}(X))\subseteq\mathrm{Non\text{-}sat}(Y)\);
\item \(f(\mathrm{Fin\text{-}sat}(X))\subseteq\mathrm{Fin\text{-}sat}(Y)\)
\end{enumerate}
\cite[Sect.~2.1.2, Def.~2.1.38]{BorgerGradelGurevich1997}.  As in
Definition~\ref{def:p02-w02-p02-036-conservative}, ``recursive function''
means a function on codes computed by a Turing machine that halts on every
member of \(X\).
\index{semi-conservative reduction}\index{reduction!semi-conservative}
\end{hbdefinition}

Every conservative reduction is semi-conservative: if \(\psi\) is
unsatisfiable, so is \(g(\psi)\) by clause (i) of
Definition~\ref{def:p02-w02-p02-036-conservative}, and if \(\psi\) has a
finite model, so has \(g(\psi)\) by clause (ii).  A semi-conservative
reduction may send an infinity axiom into any of the three classes of \(Y\),
and it may send a satisfiable formula without a finite model to an
unsatisfiable one.

\begin{hbtheorem}{thm:p02-w02b-p02-031-semi-inseparable}{proved=annex:p02-w02b-p02-031-semi-inseparable}{Semi-conservative reductions preserve inseparability}
Let \(f\) be a semi-conservative reduction from \(X\) to \(Y\).  If
\(\mathrm{Non\text{-}sat}(X)\) and \(\mathrm{Fin\text{-}sat}(X)\) are
recursively inseparable, then so are \(\mathrm{Non\text{-}sat}(Y)\) and
\(\mathrm{Fin\text{-}sat}(Y)\)
\cite[Sect.~2.1.2, Exercise~2.1.40(1)]{BorgerGradelGurevich1997}.
\end{hbtheorem}

With \(X=\mathrm{FO}\), Theorem~\ref{thm:p02-w02-p02-036-inseparable} supplies
the hypothesis.  So every recursive class \(Y\) that receives a
semi-conservative reduction from \(\mathrm{FO}\) has recursively inseparable
sets \(\mathrm{Non\text{-}sat}(Y)\) and \(\mathrm{Fin\text{-}sat}(Y)\); in
particular neither of them is decidable.

The second item of the source's Exercise~2.1.40 asks for a recursive
\(g:X\to Y\) with \(\psi\) (finitely) satisfiable if and only if \(g(\psi)\) is
(finitely) satisfiable, for every \(\psi\in X\), given a semi-conservative
reduction \(f\) from \(X\) to \(Y\)
\cite[Sect.~2.1.2, Exercise~2.1.40(2)]{BorgerGradelGurevich1997}.  Read for
arbitrary recursive classes \(X\) and \(Y\), that claim is false, as the next
lemma shows.  For \(X=\mathrm{FO}\) it is true: it is the content of the
source's Theorem~2.1.39, attributed to Gurevich, which is proved in
\cref{sec:p02-w02b-p02-031-gurevich}
(\cref{thm:p02-w02b-p02-031-gurevich,thm:p02-w02b-p02-031-ex-2-1-40-2}).

\begin{hblemma}{thm:p02-w02b-p02-031-semi-not-conservative}{proved=annex:p02-w02b-p02-031-semi-not-conservative}{A semi-conservative reduction without a conservative one}
Let \(K\) be a binary relation symbol, let \(\prec_K\) be the infinity axiom
of Example~\ref{def:p02-w02-p02-036-infinity-axiom}, and put
\(X\defeq\{\prec_K\}\) and \(Y\defeq\{\exists x\,Kxx\}\).  Both classes are
recursive.  The constant map \(f:X\to Y\) is a semi-conservative reduction
from \(X\) to \(Y\), and there is no map \(g:X\to Y\) such that for every
\(\psi\in X\), \(\psi\) is satisfiable if and only if \(g(\psi)\) is
satisfiable and \(\psi\) is finitely satisfiable if and only if \(g(\psi)\) is
finitely satisfiable.
\end{hblemma}

The example turns on a class \(Y\) without infinity axioms.  If \(X\subseteq
\mathrm{FO}\) and \(Y\) is a conservative reduction class, the second item holds
for trivial reasons: a conservative reduction \(G\) from \(\mathrm{FO}\) to
\(Y\), restricted to \(X\), is a map of the required kind, and it is recursive
on \(X\) because \(G\) halts on every member of \(\mathrm{FO}\).

\subsection{Effective inseparability and Gurevich's theorem}
\label{sec:p02-w02b-p02-031-gurevich}

The source's Theorem~2.1.39 says that, for reductions out of the class of
all formulas, the semi-conservative requirement costs nothing: whenever a
recursive class \(X\) receives a semi-conservative reduction from
\(\mathrm{FO}\), it also receives a conservative one.  The source indicates
the proof in three moves \cite[Sect.~2.1.2, pp.~37--38]{BorgerGradelGurevich1997}:
a strengthening of recursive inseparability, called effective
inseparability; the observation that \(\mathrm{Non\text{-}sat}\) and
\(\mathrm{Fin\text{-}sat}\), and with them their images in \(X\), are
effectively inseparable; and an embedding theorem of recursion theory, by
which any disjoint pair of computably enumerable sets can be pulled back
exactly from an effectively inseparable computably enumerable pair.  This
subsection states the three moves for the machine model and coding of
\cref{sec:p02-w02-p02-036-inseparability}; the annex proves each of them.

\begin{hbconvention}{def:p02-w02b-p02-031-we}{Indices of computably enumerable sets}
Let \(V\) be the decidable set of valid machine codes and \(M_e\) the machine
decoded from \(e\in V\), as in the paragraph before
Definition~\ref{def:p02-w02-p02-036-halting}.  For \(e\in\mathbb N\) put
\[
\begin{aligned}
 W_e\defeq\{n\in\mathbb N:{}& e\in V\text{ and the computation of }M_e\
   &\text{on }\operatorname{bin}(n)\text{ reaches }q_{\mathrm{acc}}\}.
\end{aligned}
\]
So \(W_e\) is the set recognized by \(M_e\), and \(W_e=\varnothing\) for an
invalid code.  The number \(e\) is called an \emph{index} of \(W_e\).
Tuples of numbers are coded by the bijective codes \(\operatorname{code}_k\)
of \cref{sec:turing-machines-01}, whose inverses are computable.
\index{index!of a computably enumerable set}
\end{hbconvention}

\begin{hbdefinition}{def:p02-w02b-p02-031-effectively-inseparable}{Effective inseparability}
Two disjoint sets \(A,B\subseteq\mathbb N\) are \emph{effectively
inseparable} if there is a total computable function
\(p:\mathbb N^{2}\to\mathbb N\) such that for all \(i,j\in\mathbb N\):
\[
 A\subseteq W_i,\quad B\subseteq W_j,\quad W_i\cap W_j=\varnothing
 \quad\Longrightarrow\quad p(i,j)\notin W_i\cup W_j
\]
\cite[Sect.~2.1.2, p.~37]{BorgerGradelGurevich1997}.  The function \(p\) is a
\emph{witness} of the effective inseparability.
\index{effective inseparability}\index{inseparability!effective}
\end{hbdefinition}

Informally, \(p\) turns any pair of disjoint computably enumerable
approximations from outside, one containing \(A\) and the other \(B\), into a
number that both of them miss.  Recursive inseparability only says that the
two approximations cannot together cover \(\mathbb N\); effective
inseparability asks for an explicit uncovered number, computed from the
indices.

\begin{hblemma}{thm:p02-w02b-p02-031-code-constructions}{proved=annex:p02-w02b-p02-031-code-constructions}{Index constructions}
\begin{enumerate}
\item Every computably enumerable set is \(W_e\) for some \(e\in V\).
\item For every total computable \(h:\mathbb N^{2}\to\mathbb N\) there is a
      total computable \(a:\mathbb N^{2}\to\mathbb N\) such that
      \(W_{a(i,c)}=\{n\in\mathbb N:h(c,n)\in W_i\}\) for all
      \(i,c\in\mathbb N\).
\item There is a total computable \(d:\mathbb N^{2}\to V\) such that for all
      \(i,j,n\in\mathbb N\): if \(n\in W_i\setminus W_j\), the computation of
      \(M_{d(i,j)}\) on \(\operatorname{bin}(n)\) reaches \(q_{\mathrm{acc}}\);
      if \(n\in W_j\setminus W_i\), it reaches \(q_{\mathrm{rej}}\); if
      \(n\notin W_i\cup W_j\), it never halts.
\item (Fixed point.)  For every computably enumerable
      \(\Theta\subseteq\mathbb N\) there is \(k\in V\) such that
      \(W_k=\{z\in\mathbb N:\operatorname{code}_2(k,z)\in\Theta\}\).
\end{enumerate}
\end{hblemma}

Item~4 is a form of Kleene's recursion theorem
(\cref{thm:recursion-theorem}) for the indices of
Convention~\ref{def:p02-w02b-p02-031-we}; the annex derives it from items 1
and~2 and does not use the recursion theorem.

\begin{hblemma}{thm:p02-w02b-p02-031-eff-transfer}{proved=annex:p02-w02b-p02-031-eff-transfer}{Elementary properties of effective inseparability}
Let \(A\) and \(B\) be effectively inseparable, with witness \(p\).
\begin{enumerate}
\item \(B\) and \(A\) are effectively inseparable, with witness
      \((i,j)\mapsto p(j,i)\).
\item \(A\) and \(B\) are recursively inseparable.
\item If \(A\subseteq A'\), \(B\subseteq B'\) and \(A'\cap B'=\varnothing\),
      then \(A'\) and \(B'\) are effectively inseparable, with witness \(p\).
\item If \(h:\mathbb N\to\mathbb N\) is total computable, \(h(A)\subseteq C\),
      \(h(B)\subseteq D\) and \(C\cap D=\varnothing\), then \(C\) and \(D\) are
      effectively inseparable.
\end{enumerate}
\end{hblemma}

\begin{hblemma}{thm:p02-w02b-p02-031-eff-halting}{proved=annex:p02-w02b-p02-031-eff-halting}{Effectively inseparable halting problems}
Put \(D_{\mathrm{rej}}\defeq\{e\in V:M_e\text{ on }\operatorname{bin}(e)
\text{ reaches }q_{\mathrm{rej}}\}\) and \(D_{\mathrm{acc}}\defeq\{e\in V:
M_e\text{ on }\operatorname{bin}(e)\text{ reaches }q_{\mathrm{acc}}\}\).  The
sets \(D_{\mathrm{rej}}\) and \(D_{\mathrm{acc}}\) are effectively
inseparable, and so are the halting problems \(H_1\) and \(H_2\) of
Definition~\ref{def:p02-w02-p02-036-halting}.
\end{hblemma}

The source asserts the second half for its halting problems and transfers it
along its reduction property \cite[Sect.~2.1.2, p.~38]{BorgerGradelGurevich1997}.

\begin{hbtheorem}{thm:p02-w02b-p02-031-eff-classes}{proved=annex:p02-w02b-p02-031-eff-classes}{Unsatisfiable and finitely satisfiable formulas are effectively inseparable}
Let \(\Lang\) be a finite relational signature containing a binary relation
symbol.  The sets \(\mathrm{Non\text{-}sat}(\operatorname{Sent}(\Lang))\) and
\(\mathrm{Fin\text{-}sat}(\operatorname{Sent}(\Lang))\) are effectively
inseparable, and so are \(\mathrm{Non\text{-}sat}\) and
\(\mathrm{Fin\text{-}sat}\).  The sets \(\mathrm{Non\text{-}sat}\) and
\(\mathrm{Fin\text{-}sat}\) are computably enumerable.  If \(X\subseteq
\mathrm{FO}\) is recursive and \(f\) is a semi-conservative reduction from
\(\mathrm{FO}\) to \(X\), then \(\mathrm{Non\text{-}sat}(X)\) and
\(\mathrm{Fin\text{-}sat}(X)\) are computably enumerable and effectively
inseparable.
\end{hbtheorem}

\begin{hbtheorem}{thm:p02-w02b-p02-031-embedding}{proved=annex:p02-w02b-p02-031-embedding}{Embedding into an effectively inseparable pair}
Let \(P_1,P_2\subseteq\mathbb N\) be disjoint computably enumerable sets, and
let \(R_1,R_2\subseteq\mathbb N\) be disjoint, computably enumerable and
effectively inseparable.  Then there is a total computable
\(g:\mathbb N\to\mathbb N\) with
\[
 P_1=g^{-1}(R_1)\qquad\text{and}\qquad P_2=g^{-1}(R_2).
\]
\end{hbtheorem}

This is the recursion-theoretic argument the source invokes for
Theorem~2.1.39 \cite[Sect.~2.1.2, p.~37]{BorgerGradelGurevich1997}.  The
map \(g\) does more than a many-one reduction of \(P_1\) to \(R_1\): it sends
every number outside \(P_1\cup P_2\) outside \(R_1\cup R_2\).

\begin{hbtheorem}{thm:p02-w02b-p02-031-gurevich}{proved=annex:p02-w02b-p02-031-gurevich}{Gurevich: semi-conservative reductions from FO suffice}
If there exists a semi-conservative reduction from the set \(\mathrm{FO}\) of
all first-order formulas to a recursive class \(X\subseteq\mathrm{FO}\), then
\(X\) is a conservative reduction class
\cite[Sect.~2.1.2, Thm.~2.1.39]{BorgerGradelGurevich1997}.
\end{hbtheorem}

\begin{hbcorollary}{thm:p02-w02b-p02-031-ex-2-1-40-2}{proved=annex:p02-w02b-p02-031-ex-2-1-40-2}{The second item of Exercise 2.1.40 for \(X=\mathrm{FO}\)}
Let \(Y\subseteq\mathrm{FO}\) be a recursive class and \(f\) a
semi-conservative reduction from \(\mathrm{FO}\) to \(Y\).  Then there is a
recursive function \(g:\mathrm{FO}\to Y\) such that every \(\psi\in
\mathrm{FO}\) is satisfiable if and only if \(g(\psi)\) is satisfiable, and
finitely satisfiable if and only if \(g(\psi)\) is finitely satisfiable
\cite[Sect.~2.1.2, Exercise~2.1.40(2)]{BorgerGradelGurevich1997}.
\end{hbcorollary}

\begin{hbexample}{def:p02-w02b-p02-031-gurevich-example}{What the conservative reduction does with an infinity axiom}
Let \(f\) be a semi-conservative reduction from \(\mathrm{FO}\) to a recursive
\(X\), and let \(\prec_K\) be the infinity axiom of
Example~\ref{def:p02-w02-p02-036-infinity-axiom}.  Nothing in
Definition~\ref{def:p02-w02b-p02-031-semi-conservative} constrains
\(f(\prec_K)\): it may be unsatisfiable, finitely satisfiable, or an infinity
axiom.  The conservative reduction \(G\) built in the proof of
\cref{thm:p02-w02b-p02-031-gurevich} sends \(\prec_K\) to a member of
\(\mathrm{Inf\text{-}axioms}(X)\).  So every recursive class that receives a
semi-conservative reduction from \(\mathrm{FO}\) contains an infinity axiom.
The construction of \(G\) gives no description of which formula of \(X\) this
is: \(G(\prec_K)\) is the number \(p(i,j)\) produced by the witness \(p\) of
effective inseparability for indices \(i,j\) obtained from a fixed point.
\end{hbexample}

\begin{hbwarning}{def:p02-w02b-p02-031-gurevich-limits}{Limits of Gurevich's theorem}
\begin{enumerate}
\item The theorem concerns reductions out of \(\mathrm{FO}\).  For an
  arbitrary recursive domain class in place of \(\mathrm{FO}\) the
  corresponding conclusion fails
  (\cref{thm:p02-w02b-p02-031-semi-not-conservative}).
\item The conservative reduction \(G\) is computable, but it is obtained from
  a fixed point (\cref{thm:p02-w02b-p02-031-code-constructions}(4)); it does
  not preserve the syntactic shape of formulas and gives no bound on the size
  of finite models.
\item The proof uses that \(\mathrm{Non\text{-}sat}(X)\) and
  \(\mathrm{Fin\text{-}sat}(X)\) are computably enumerable, which holds
  because \(X\) is recursive.
\end{enumerate}
\end{hbwarning}

\emph{Relationships.}  \Cref{thm:p02-w02b-p02-031-register-reduction} is the
register-machine counterpart of \cref{thm:p02-w02-p02-036-refined}: there the
time and the tape cells are elements of the model, here the register contents
are, and the time disappears.  \Cref{thm:p02-w02b-p02-031-ex-2-1-37} answers,
for two Krom--Horn prefix classes, the question which classes are conservative
reduction classes in the sense of
Definition~\ref{def:p02-w02-p02-036-conservative}; with
\cref{thm:p02-w02-p02-036-conservative} it gives pairwise recursive
inseparability of the three countermodel classes inside each of them.  The
logic machine of \cref{thm:p02-w02b-p02-031-logic-machine} is built from the
index constructions of \cref{sec:p02-w02b-p02-031-gurevich}.
\Cref{thm:p02-w02b-p02-031-val-to-finsat} reuses the
run sentence, the path translation and the renaming of
\cref{sec:p02-w02-p02-036-finite-validity} and the enumerability of validity
(\cref{thm:fol-validity-enumerable}); it complements
\cref{thm:p02-w02-p02-036-trakhtenbrot}, which reduces the self-application set
\(K\) of \cref{sec:universal-reductions-03} to finite satisfiability.
\Cref{thm:p02-w02b-p02-031-semi-inseparable} is the semi-conservative analogue
of \cref{thm:p02-w02-p02-036-conservative}, with one pair of classes in place
of three.  \Cref{thm:p02-w02b-p02-031-eff-halting} strengthens
Lemma~\ref{thm:p02-w02-p02-036-halting} for the pair \(H_1\), \(H_2\), and
\cref{thm:p02-w02b-p02-031-eff-classes} strengthens the pair
\(\mathrm{Non\text{-}sat}\), \(\mathrm{Fin\text{-}sat}\) of
\cref{thm:p02-w02-p02-036-inseparable}.  The fixed-point item of
\cref{thm:p02-w02b-p02-031-code-constructions} plays the part of Kleene's
recursion theorem (\cref{thm:recursion-theorem}), and the index
constructions use the parameter idea of \cref{thm:parameter} for machine
codes.

\emph{Sources.}  Definition~2.1.14 (p.~27), Theorem~2.1.15 and
Corollary~2.1.16 with their formulas (pp.~28--29) and Exercise~2.1.17
(p.~29) of \cite[Sect.~2.1.1]{BorgerGradelGurevich1997}; Exercises 2.1.32 and
2.1.34 with the refined formulas (p.~36) and Exercise~2.1.37 (p.~37) of
\cite[Sect.~2.1.2]{BorgerGradelGurevich1997}; the prefix-vocabulary classes of
\cite[Sect.~1.3, Def.~1.3.1, p.~8]{BorgerGradelGurevich1997}.  The simulation
of Turing machines by two-register machines, which the source takes from the
literature, is proved in the annex for this handbook's machine model.
Exercise~2.1.33 (p.~36), Definition~2.1.38, Theorem~2.1.39
with its proof outline (pp.~37--38) and Exercise~2.1.40 (p.~38) of
\cite[Sect.~2.1.2]{BorgerGradelGurevich1997}.  The source names effective
inseparability and the embedding theorem and says how they combine; the
proofs of \cref{thm:p02-w02b-p02-031-code-constructions,thm:p02-w02b-p02-031-eff-transfer,thm:p02-w02b-p02-031-eff-halting,thm:p02-w02b-p02-031-eff-classes,thm:p02-w02b-p02-031-embedding}
in the annex are written for this handbook's machine model and coding; the proof of
\cref{thm:p02-w02b-p02-031-val-to-finsat} does not follow the source's hint
through Herbrand's theorem.
